Este trabalho apresenta a implementação, análise e comparação de diferentes algoritmos de agrupamento aplicados aos conjuntos de dados Adult e Dry Bean. Foram avaliados os métodos KMeans, Fuzzy C-Means, Gustafson-Kessel e Agglomerativo, com o objetivo de investigar sua capacidade de identificar padrões e separar grupos em dados reais e multivariados. A metodologia envolveu o pré-processamento dos dados, aplicação dos algoritmos, avaliação dos resultados por métricas quantitativas (acurácia, precisão, revocação, especificidade, F1-score e matriz de confusão) e visualização dos agrupamentos por meio de PCA em duas dimensões. Os resultados obtidos evidenciam as características, vantagens e limitações de cada abordagem, destacando a importância da escolha do método de agrupamento conforme a natureza dos dados e o objetivo da análise. O estudo contribui de forma didática para a compreensão dos principais algoritmos de clustering e suas aplicações práticas em ciência de dados.